\section{Literature Review}
\label{sec:literature-review}

% STATUS: substantially expanded, still not final. references.bib now has
% 15 settled/foundational + 14 recent (2023-2025) citations = 29 total,
% comfortably clearing publication_supplement.md's >=20-reference checklist
% item; 14 of 29 are 2023-2025, just short of its >=15-recent target. Every
% citation below (recent and foundational alike) has been verified against
% Crossref metadata (see references.bib comments) -- do not add a citation
% here without the same verification. Two citation flags remain genuinely
% open (marked TODO below) -- resolve both before treating this section as
% submission-ready.

This review is organized by the same four MPPT design paradigms compared
in this paper, both to survey prior work and to make the cross-paradigm
framing of Section~\ref{sec:introduction} concrete, followed by recent
cross-paradigm reviews that position this paper's own contribution.

\subsection{Perturbative Methods}
Perturb-and-observe remains the most widely deployed MPPT method for its
simplicity and low computational cost; \cite{femia2005optimization}
provides the adaptive-step-size design this paper's P\&O implementation
follows, trading off oscillation amplitude against tracking speed by
scaling the perturbation step with $|dP/dV|$. Incremental conductance
\cite{hussein1995maximum} refines this by directly comparing $dI/dV$ against
$-I/V$, in principle eliminating steady-state oscillation around the true
maximum power point (MPP) at the cost of a boundary-condition singularity
as $dV \to 0$ that must be handled explicitly. Adaptive step-size tuning
for both methods remains an active research area:
\cite{ibrahim2023optimizing} uses particle swarm optimization to tune the
step size of both P\&O and IncCond for a grid-tied system, reporting
improved steady-state oscillation over fixed-step baselines -- the same
tradeoff this paper's adaptive-step P\&O design targets, though via a
simpler analytical rule rather than a metaheuristic-tuned one.

\subsection{Rule-Based Methods}
Fuzzy logic control removes the need for an explicit analytical model of
the PV/converter system by encoding tracking behavior as a linguistic rule
base. \cite{zhu2021comprehensive} surveys four FLC-MPPT input/output
variable choices ($dP\&dV$, $dP\&dI$, $E\&dE$, $G\&F$) at a $5\times5$
partition size; this paper's design instead uses a $7\times7$ partition
following common (E, $\Delta E$) $\to \Delta D$ conventions in the wider
fuzzy-MPPT literature. \cite{sunkara2024novel} presents a hybrid fuzzy
logic MPPT controller specifically evaluated under partial shading,
reporting improved global-search behavior over a standalone fuzzy
controller -- directly relevant context for this paper's own finding
(Section~\ref{sec:results}) that a non-hybrid fuzzy controller, like every
other non-hybrid algorithm tested, fails to escape a local maximum under
partial shading; a hybrid extension is exactly the kind of change that
finding motivates as future work.
% TODO: the specific 7x7 partition-size literature precedent (as opposed to
% 5x5, which \cite{zhu2021comprehensive} surveys directly) is still an open
% citation flag per CLAUDE.md -- find and verify a paper using the same
% 7x7 partition before final submission, rather than citing this 5x5 survey
% as if it establishes 7x7 precedent.

\subsection{Learning-Based Methods}
\cite{kofinas2017reinforcement} is the foundational reference for
model-free reinforcement learning applied to MPPT, using tabular Q-learning
over a discretized (voltage, $dP/dV$) state space -- the same state
representation this paper's Q-learning implementation uses.
\cite{chen2023research} studies Q-table design specifically, proposing an
improved update scheme for the same tabular setting this paper uses rather
than moving to function approximation or deep RL -- a natural point of
comparison for whether this paper's Q-learning generalization failure
(Section~\ref{sec:partial-shading-results}) is a consequence of the
particular Q-table design or a more fundamental limitation of tabular
discretization under distribution shift. A key open question for any
learning-based MPPT controller is whether it generalizes to conditions
outside its training distribution or merely memorizes them;
\cite{phan2020deep} studies a deep-RL controller specifically under partial
shading, a harder and more directly relevant generalization test than
flat-irradiance held-out evaluation alone, and is a natural point of
comparison for this paper's own partial-shading generalization finding
(Section~\ref{sec:results}), where the Q-learning agent -- trained only on
flat, single-module irradiance -- fails on multi-module shading in a
qualitatively different way than the perturbative and model-based
algorithms.

\subsection{Model-Based Nonlinear Control}
Sliding mode control (SMC) applies variable-structure control theory
\cite{utkin1977variable} to MPPT by defining a sliding surface (commonly
$s = dP/dV$, driving the system toward the MPP condition) and a reaching
law governing convergence to that surface. This paper's exponential
reaching law follows \cite{gao1993variable}, and its boundary-layer
chattering mitigation -- replacing $\mathrm{sign}(s)$ with a saturation
function -- follows the classical treatment in \cite{slotine1991applied}.
Chattering mitigation remains an active SMC-MPPT research area:
\cite{liu2025photovoltaic} combines PID-type SMC with an improved grey
wolf optimizer to tune controller gains rather than relying on a
hand-tuned boundary layer as this paper does, illustrating one direction
(automated gain tuning) this paper's grid-searched
\texttt{SMC\_K}/\texttt{SMC\_Q}/\texttt{SMC\_BOUNDARY} constants
(\texttt{config.py}) could be extended toward.
% TODO: per CLAUDE.md's own caution, no PV-specific SMC citation was
% finalized before implementation; \cite{zhang2022novel},
% \cite{mostafa2020tracking}, and \cite{contreras2024novel} are
% Crossref-verified candidates found during this literature search, but
% none has been confirmed to use the exact same reaching-law family as this
% implementation (zhang2022novel uses a multi-power reaching law + sigmoid,
% not confirmed exponential + boundary layer; mostafa2020tracking's and
% contreras2024novel's content weren't independently verifiable via
% automated fetch). Read all three in full before citing any as the
% specific design match.

\subsection{PV Modeling and Boost Converter Design}
The boost converter sizing and small-signal analysis in this paper follow
the standard treatment in \cite{erickson2001fundamentals}, including the
right-half-plane zero in the control-to-output transfer function that
simplified worked examples sometimes omit. \cite{serra2024control}
addresses boost-converter control specifically in discontinuous conduction
mode (DCM) under a constant-power load -- directly relevant to this
paper's own DCM finding (Section~\ref{sec:converter-model}), that this
design's converter enters DCM at light load and that every algorithm's
duty-cycle-perturbation reasoning implicitly assumes the CCM small-signal
model; \cite{serra2024control}'s dedicated DCM control strategy is one
concrete way a future revision could remove that confound rather than
merely flag it as a limitation.

\cite{villalva2009comprehensive} and \cite{ishaque2011simple} both address
PV array modeling but at different levels: the former's contribution is a
single-diode, three-datasheet-point parameter extraction method, while the
two-diode model used in this paper -- which better represents recombination
currents at low voltage and is important for correctly resolving multiple
local maxima under partial shading -- follows \cite{ishaque2011simple}
specifically; the two should not be conflated when citing the two-diode
model's origin. More recent parameter-extraction work continues to refine
this: \cite{syed2024parameter} proposes an analytical (non-iterative)
two-diode extraction method, and \cite{moshksar2024two} compares two
datasheet-based double-diode extraction approaches on an accuracy/
simplicity tradeoff -- directly relevant to this paper's own extraction
method (Section~\ref{sec:parameter-extraction}), which is explicitly
flagged as a simplification (fixed $a_1$, analytically-solved
$I_{s1}=I_{s2}$) relative to the full iterative procedure in
\cite{ishaque2011simple}; both recent papers are natural points of
comparison for tightening that simplification in future work.
\cite{abad2025analytical} provides a recent analytical treatment of
bypass- and blocking-diode effects on the partial-shading $P$-$V$ curve,
complementing this paper's simulation-based bypass-diode composition
(Section~\ref{sec:pv-model}) with an independent analytical derivation of
the same multiple-local-maxima phenomenon this paper's Monte Carlo results
depend on.

\subsection{Cross-Paradigm Positioning}
The most direct historical point of comparison is
\cite{esram2007comparison}, whose 19-technique survey remains, two decades
later, the most widely-cited broad MPPT comparison in this journal. Three
things distinguish this paper from it, beyond the two decades of algorithm
development in between: \cite{esram2007comparison} is a narrative survey of
results reported by different authors under different conditions, not a
common statistically-validated experimental comparison; it predates
learning-based and model-based nonlinear control as MPPT paradigms
entirely, so a four-paradigm framing was not yet possible; and it does not
release runnable code, so its comparisons cannot be independently
re-executed the way every result in this paper can be.

\cite{endiz2025review} and \cite{yaqoob2025advanced} are both recent,
broad reviews spanning classical, metaheuristic, and AI-based MPPT
methods -- the same cross-paradigm scope this paper argues is missing from
most single-comparison studies (Section~\ref{sec:introduction}). Neither,
however, provides the specific combination this paper does: a validated
physical model, open-source reproducible code, and Monte Carlo statistical
rigor applied uniformly across all paradigms compared, rather than a
literature synthesis of results reported under differing conditions by
different authors. \cite{abdelmalek2024experimental} and
\cite{taha2025design} are both recent examples of the "yet another
bio-inspired/hybrid metaheuristic optimizer" pattern this paper's
introduction explicitly declines to add to (a zebra optimization algorithm
and a GWO-PSO hybrid, respectively) -- included here not as directly
comparable baselines but as evidence the pattern remains active and
crowded, reinforcing rather than undermining this paper's choice to compete
on reproducibility and paradigm coverage instead. \cite{ashraf2025python}
is a recent, directly relevant example of the open-source reproducibility
movement this paper is also part of, releasing a Python framework for
metaheuristic MPPT techniques evaluated under partial shading -- but,
consistent with the pattern above, it is scoped to a single paradigm
(metaheuristic optimization) rather than the four-paradigm comparison this
paper provides.

% Remaining work for this section: one more recent citation to comfortably
% clear publication_supplement.md's >=15-recent-references target
% (currently 14 of 29 total references are 2023-2025), and resolving both
% TODOs above (7x7 fuzzy partition precedent; specific PV-SMC design-match
% citation) before this section is submission-ready.
